% Zone „Das kennst du schon“ – aus bank/binomische-formeln/zone.jsonl (zusammenbau v0.1)
\einheitenkopf[zone][Kennst du schon]{Das kennst du schon}
\setcounter{aufgabe}{0}

%% TODO Titel: Ich-kann-Satz fehlt in der Bank (Platzhalter: Kettenname) – Nr. 1
%% TODO Anweisung: ein Satz über der ersten Teilaufgabe fehlt in der Bank – Nr. 1
\begin{aufgabe}{Multiplizieren mit Vorzeichen und Punkt vor Strich}
\begin{teile}
\teil $(-4) \cdot (-8)$ – Ergebnis? \leerfeld
\teil $(-3) \cdot 5 - 2 \cdot (-7)$ – Ergebnis? \leerfeld
\end{teile}
\end{aufgabe}

%% TODO Titel: Ich-kann-Satz fehlt in der Bank (Platzhalter: Kettenname) – Nr. 2
%% TODO Anweisung: ein Satz über der ersten Teilaufgabe fehlt in der Bank – Nr. 2
\begin{aufgabe}{Termwert berechnen, auch mit negativer Zahl, als Probe einer Umformung}
\begin{teile}
\teil $3x + 5$ für $x = 4$ – Termwert? \leerfeld
\teil $2x^2 - 3x$ für $x = -2$ – Termwert? \leerfeld
\end{teile}
\end{aufgabe}

%% TODO Titel: Ich-kann-Satz fehlt in der Bank (Platzhalter: Kettenname) – Nr. 3
%% TODO Anweisung: ein Satz über der ersten Teilaufgabe fehlt in der Bank – Nr. 3
\begin{aufgabe}{Gleichartige Glieder mit einer Variablen zusammenfassen, Vorzahl eins beachten, Potenz getrennt halten}
\begin{teile}
\teil $4x + 3x$ – zusammengefasst? \leerfeld
\teil $5x - 8x + 2$ – zusammengefasst? \leerfeld
\end{teile}
\end{aufgabe}

%% TODO Titel: Ich-kann-Satz fehlt in der Bank (Platzhalter: Kettenname) – Nr. 4
%% TODO Anweisung: ein Satz über der ersten Teilaufgabe fehlt in der Bank – Nr. 4
\begin{aufgabe}{Zahl mal Klammer und Minusklammer (Distributivgesetz, alle Vorzeichen drehen)}
\begin{teile}
\teil $4 \cdot (x + 2)$ – ohne Klammer? \leerfeld
\teil $-2 \cdot (x - 6)$ – ohne Klammer? \leerfeld
\end{teile}
\end{aufgabe}

%% TODO Titel: Ich-kann-Satz fehlt in der Bank (Platzhalter: Kettenname) – Nr. 5
%% TODO Anweisung: ein Satz über der ersten Teilaufgabe fehlt in der Bank – Nr. 5
\begin{aufgabe}{Quadratzahlen bis 15² erkennen und Quadratwurzeln daraus}
\begin{teile}
\teil $\sqrt{64}$ – Wert? \leerfeld
\teil $(-6)^2$ – Wert? \leerfeld
\end{teile}
\end{aufgabe}

%% TODO Titel: Ich-kann-Satz fehlt in der Bank (Platzhalter: Kettenname) – Nr. 6
%% TODO Anweisung: ein Satz über der ersten Teilaufgabe fehlt in der Bank – Nr. 6
\begin{aufgabe}{Scheitelpunktform lesen}
\begin{teile}
\teil $y = (x - 4)^2 + 1$ – Scheitel? S( \leerfeld | \leerfeld )
\teil $y = (x - 6)^2 - 7$ – Scheitel? S( \leerfeld | \leerfeld )
\end{teile}
\end{aufgabe}

%% TODO Titel: Ich-kann-Satz fehlt in der Bank (Platzhalter: Kettenname) – Nr. 7
%% TODO Anweisung: ein Satz über der ersten Teilaufgabe fehlt in der Bank – Nr. 7
\begin{aufgabe}{Ausklammern eines gemeinsamen Zahlfaktors oder einer Variablen}
\begin{teile}
\teil $5x + 20$ – ausgeklammert? \leerfeld
\teil $4x^2 - 8x$ – ausgeklammert? \leerfeld
\end{teile}
\end{aufgabe}

%% TODO Titel: Ich-kann-Satz fehlt in der Bank (Platzhalter: Kettenname) – Nr. 8
%% TODO Anweisung: ein Satz über der ersten Teilaufgabe fehlt in der Bank – Nr. 8
\begin{aufgabe}{Zahl mal Klammer und Minusklammer (Distributivgesetz, alle Vorzeichen drehen) – Fehler finden}
\begin{teile}
\teil Finde den Fehler und rechne richtig: \rechnung{-(x - 9) &= -x - 9}
\end{teile}
\end{aufgabe}

%% TODO Titel: Ich-kann-Satz fehlt in der Bank (Platzhalter: Kettenname) – Nr. 9
%% TODO Anweisung: ein Satz über der ersten Teilaufgabe fehlt in der Bank – Nr. 9
\begin{aufgabe}{Zahl mal Klammer und Minusklammer (Distributivgesetz, alle Vorzeichen drehen) – selbst rechnen}
\begin{teile}
\teil $-(2x - 5)$ – ohne Klammer? \leerfeld
\end{teile}
\end{aufgabe}
